
\section{Mean-field analysis}
\label{sec:meanfield}

The macroscopic behaviour of the agent-based model is governed by two slow variables: the population size $N(t)$
and the mean social state $S(t)=N^{-1}\sum_i s_i$. A hierarchy of mean-field descriptions explains the regimes of
the core model without the rules, extracts the mechanism and the thresholds of the irreversibility rules R1 and R2, and gives the
quantitative predictions tested in Section~\ref{sec:res-Q}. The analysis itself is in
Appendices~\ref{app:meanfield} and~\ref{app:mf-supp}; here we state its scope and its predictions, with their
outcome.

\subsection{Scope}
\label{sec:mf-scope}

Most of the mean-field analysis describes the \emph{short-lived control} and the long-lived control (the core
model without the four rules: $a_{\max}=60$ or 120, $\alpha=1.4$, $r=0.01$, $\gamma=0.1$;
Section~\ref{sec:model-candidates}), not the calibrated candidates. In that model
reproduction is limited by mate encounter, which gives a built-in Allee effect, and deterioration is driven by the
Poisson tail of the local occupancy, which makes the social state a switch between recovery and loss at a
population $N^*$; a two-variable mean field then has a superstable limit cycle, the oscillations of the spatial
model are sustained by \emph{healthy pockets} of uncrowded individuals, and extinction is a mixture of an early
deterministic-like mode and a rare late mode, without secular drift of the peaks. Its correspondence with
\textsf{C1} is limited: the switch point ($N^*\approx6200$ with Poisson occupancy, 3900--4600 with negative-binomial
occupancy of dispersion $\xi=1$--2, 5300 for $\xi=5$) lies close to the damage-threshold capacity $K_\gamma=7200$ and
far above the peak of \textsf{C1}
($\approx2600$), which is set by aggregation and refuges; under R1 adults do not recover, so the limit cycle does not
exist in \textsf{C1}; and the first-moment closure fails where it matters (at $N=1600$ the effective weekly loss of
social state is 0.75, against 0.001 predicted with Poisson occupancy). What carries over to \textsf{C1} is the demographic law of phase D, Eq.~\eqref{eq:phaseD}, which uses only the
mortality; the capacity bound $\Phi_{\rm BO}^{\rm pers}\le\phi_R$ of Section~\ref{sec:model-groups}; and two
heuristic statements about R2 with the parameters of \textsf{C1}, the elasticity $\theta$ of Q6 and the
generational map of Q7. Under R1 recovery becomes a \emph{lineage} process, whose bottleneck is the fecundity
$s_f^2s_m$ of deteriorated parents (Table~\ref{tab:lineage}); under R2 the recovery time grows with the
social-learning weight $\lambda$ and is a lower bound for the irreversibility threshold (Eq.~\eqref{eq:trec}). A
kinetic mean field with persistent cell disorder confirms both rules in the long-lived scenario for the two rates
it was run with, $r\in\{0.01,0.03\}$: R1 gives certain extinction for every window $a_w\le24$, and R2 alone needs a
weight $\lambda$ above a threshold that grows with $r$ (defining the threshold as the first value of the grid
$\{0,0.3,0.5,0.7,0.9,1\}$ at which at least half of the 10 seeds go extinct: $\lambda_c\le0.3$ at $r=0.01$, with
8/10 seeds already at 0.3, and $\lambda_c\in(0.7,0.9]$ at $r=0.03$). These are bounds and orders of magnitude, not
predictions for \textsf{C1}.

\subsection{Predictions for the calibrated candidate and their outcome}
\label{sec:mf-predictions}

Before the production sweeps, the analysis was turned into quantitative predictions for \textsf{C1} (Q1--Q7,
Section~\ref{sec:stats}). Q2 (necessity of each rule) comes from the calibration; the other six have a mean-field
origin, and only Q3--Q4, Q6 and Q7 rest on pieces that carry over to \textsf{C1}. The outcomes below are those of
the production runs with the primary evaluator; where the preregistered evaluator gives a different verdict it is stated
(Section~\ref{sec:res-Q}).

\begin{description}
\item[Q1, transplant.] Under R1, phase-D animals ($s\approx0$) cannot found a colony, because their recovery must
  pass through reproduction (Table~\ref{tab:lineage}); with individual recovery at $r=0.1$ the time to viability,
  about 3 weeks, is shorter than the fertile time left, so they can, up to an age ceiling.
  \emph{Outcome:} the null hypothesis is rejected (0/60 against 51/60 at $1.75\,t'_{\rm pk}$) and the predicted rate
  in the enclosure without R1 and R2 ($\ge0.25$) is reached (0.85); at the later sampling time of the preregistered
  protocol the same animals give 12/60, because of the age ceiling. Since the transplanted animals have $s\approx0$
  and fecundity scales as $s_f^2s_m$, the 0/60 follows from R1 by construction: Q1 is a consistency check of R1, as
  Q4 is of the demography. Transplants of intermediate animals confirm the bottleneck (8\% found a colony under R1
  against 95\% without it) but not the collapse of the lineage: their offspring settle at the parents' level.
\item[Q3, Q4, the phase-D law.] Once births stop, extinction is the death of the last survivor of a closed, ageing
  population. With $n$ animals born at the last birth event,
  \begin{equation}
    \mathbb E[T_{\rm ext}-t_{\rm last}]\simeq\sum_{a\ge1}\Bigl[1-\bigl(1-S(a)\bigr)^n\Bigr],\qquad
    S(a)=\prod_{b=1}^{a}\bigl[1-h(b)\bigr],\qquad
    h(b)=\bigl[1+e^{-k(b-a_{\max})}\bigr]^{-1},
    \label{eq:phaseD}
  \end{equation}
  where $S(a)$ is the probability that an animal born at $t_{\rm last}$ is still alive $a-1$ weeks later (the
  product starts at $b=1$ because the first mortality draw happens at age one, in the week of birth). The
  expectation grows one-to-one with $a_{\max}$ and only logarithmically with $n$.
  \emph{Outcome:} confirmed (slope $0.996\pm0.014$); the logarithmic dependence on $n$ appears in the lattice-size
  stage. Since this follows from the demography alone, it is a consistency check rather than a test of the rules.
\item[Q5, preferred occupancy.] A deteriorated individual weighs a cell with $n$ acquaintances by
  $(1+\alpha n)e^{-\kappa n}$, maximal at $n^*=1/\kappa-1/\alpha$ (Eq.~\eqref{eq:model-nstar}); pools above $m_c$
  and isolated animals were predicted to coexist only if $n^*\approx m_c$. The prediction was derived and calibrated
  on \textsf{C0}, which has no refuges, and the preregistration should have anticipated that R7 decouples the
  isolated class from this balance. \emph{Outcome:} refuted in \textsf{C1}.
\item[Q6, window against rate.] Under R2 with $\lambda=1$ a juvenile ends the window with
  $s(a_w)=w\,s_f+r\int_0^{a_w}\bar s_V(a)\,da$; the elasticity of the learned part is one with respect to $r$ and
  $\theta=\bar s_V(a_w)/\langle\bar s_V\rangle_{[0,a_w]}<1$ with respect to $a_w$, if the competence of the
  neighbours decreases with age.
  \emph{Outcome:} the prediction $0<\theta<1$ is not supported with the primary evaluator: the Calhoun-like boundary
  of the plane is set by the learning rate and the window has little or no effect ($\hat\theta=-0.15$
  [$-0.20$, $-0.05$] in a single-index model that an additive model beats, $\Delta{\rm QAIC}=28$); with the
  preregistered evaluator the single-index model is the preferred one and gives $\hat\theta=0.30$ [0.15, 0.50], as in
  the runs of the preregistered design. The difference is the criterion O11: measured on the matured late cohorts, the boundary is set by their
  socialisation at fast learning, where a longer window fails less often, so the weak dependence on $a_w$ runs
  against the predicted sign without amounting to a reversal. The elasticity argument describes the learned
  competence at the end of the window, not which cohorts end up below the threshold.
\item[Q7, generational map.] Without crowding, the mean competence at maturity obeys
  \begin{equation}
    S_{g+1}=w\,\tilde S_g+\varphi\,\Lambda\,S_g,\qquad \Lambda=r\,a_w,\qquad
    \tilde S_g=\frac{\mathbb E[s^3]}{\mathbb E[s^2]}\equiv\sigma S_g,
    \label{eq:genmap}
  \end{equation}
  where $\sigma\ge1$ is the selection factor of the birth flux and $\varphi$ the competence of a pup's neighbours
  relative to $S$. Competence dies out if $\sigma w+\varphi\Lambda<1$, and since the reproduction number
  $R_0(S)\approx70\,S^3$ of the long-lived demography falls below one for $S<0.24$, the colony then goes extinct
  \emph{without} crowding; at the \textsf{C1} values ($w=0.2$, $\Lambda=1.2$) the map is supercritical.
  \emph{Outcome:} the test confirms an endogenous collapse without crowding in the subcritical corner, which is not
  Calhoun-like because the colonies first overshoot the capacity. With $a_w$ varied, the transition in
  $\Lambda=r\,a_w$ falls in the same grid bracket for $a_w=12$ and 24 and one grid step lower for $a_w=6$, so the
  exponent of $a_w$ relative to $r$ lies in $(0,1)$ between 6 and 12 weeks and is compatible with 1 between 12 and
  24; the grid resolves neither the preregistered equivalence margin $[2/3,3/2]$ nor a single exponent (51 of 54
  points have $P_{\rm ext}\in\{0,1\}$), so the product form holds only approximately and short windows are less
  effective than it predicts. The linearity of the boundary in $(w,\Lambda)$ is imposed by the logistic predictor,
  and the fitted $\hat\varphi\approx1.1$--1.2, above the preregistered $\varphi\le1$, has no valid interval.
\end{description}
The status of the exploratory mean-field statements is listed in Appendix~\ref{app:mf-supp}.
